Let $\mathcal S$ be a family of \emph{finite} (simple, undirected) graphs.
Assume that for every integer $n\ge 1$ there exists a graph $G_n\in\mathcal S$ such that
\[
  G_n \to (K_3)^{e}_n,
\]
meaning: for every edge-colouring $c:E(G_n)\to\{1,\dots,n\}$ there is a monochromatic triangle.

Question (as written). Must it then be true that for every infinite cardinal $\kappa$ there exists a graph $G$ such that

\begin{enumerate}[label=(\alph*)]
\item every finite subgraph of $G$ belongs to $\mathcal S$, and
\item $G\to (K_3)^{e}_\kappa$ (every $\kappa$-edge-colouring of $G$ contains a monochromatic triangle)?
\end{enumerate}
